\chapter{Reading a search tree like a chess player}
\label{ch:readingtree}

\begin{objectives}
\begin{itemize}
\item Read a verbose search dump in a fixed order and know what each number licenses you to
  say.
\item Compute the divergence between policy and visits, and classify a position as
  \emph{agreement}, \emph{optical trap} or \emph{hidden gem} using real output.
\item Measure how much a search added to the network's static opinion, in units you can
  quote.
\item Decide, from $n$ alone, how many decimal places of an evaluation are meaningful.
\item Write down what a training tool may honestly say about a position, and what it may not.
\end{itemize}
\end{objectives}

\section{The order to read them in}

A verbose dump has too many numbers. Read them in this order; each answers a different
question, and each depends on the ones above it.

\begin{center}
\begin{tabular}{@{}rlp{8.3cm}@{}}
\toprule
 & Quantity & The question it answers \\
\midrule
1 & $N$ of the root & \emph{How much thinking is this analysis based on?} Everything below is
  conditioned on this. \\
2 & $d$ at the root & \emph{Is anything at stake?} (Chapter~\ref{ch:currency}.) \\
3 & the visit distribution & \emph{How did the engine allocate belief across plans?} \\
4 & $Q$ of the top move & \emph{Who stands better, after search?} \\
5 & $Q$ versus $V$ of the top move & \emph{How much did calculation change the picture?} \\
6 & $P$ versus visit share & \emph{Where did intuition and calculation disagree?} \\
7 & the principal variation & \emph{What is the plan?} --- read last, and trust least. \\
\bottomrule
\end{tabular}
\end{center}

\begin{pitfall}
Most people read the PV first, because it is the most chess-like object on the screen. It is
also the least reliable: the PV beyond the first two or three plies is often a thin line
visited a handful of times. Chapter~\ref{ch:onemachine}'s $1/\sqrt{n}$ applies with full
force --- the tenth move of a PV may rest on three samples. Read the visit distribution first;
it is the part backed by the whole budget.
\end{pitfall}

\section{The visit distribution is the engine's opinion}

Chapter~\ref{ch:puct} proved that visits reproduce the prior when search learns nothing, and
sharpen it when search learns something. So the visit distribution is not a technical
detail --- it is \emph{the output}. It is the engine's answer to "how seriously should each
plan be taken?", and it is the closest thing an engine produces to a human's sense of which
candidate moves are live.

Three shapes recur, and they mean different things:

\begin{center}
\begin{tabular}{@{}lll@{}}
\toprule
Shape & Example & Reading \\
\midrule
Concentrated & $90/4/3/2/1$ & one move; the rest are refuted or clearly worse \\
Split & $38/32/17/7/5$ & a genuine choice; several playable plans \\
Flat & $21/21/20/19/19$ & nothing is at stake, or nothing has been learned \\
\bottomrule
\end{tabular}
\end{center}

A flat distribution needs care: it means one of two opposite things --- the position is dead
(Chapter~\ref{ch:puct}'s drawn ending, where all five moves drew and visits stayed within
$3\%$ of each other) or the search was too short to distinguish anything. Item 1 of the
reading order settles which.

\begin{worked}{quantifying the shape}
The natural summary is the \textbf{entropy} of the visit distribution
$\pi(a) = N(s,a)/N(s)$:
\[
  \Ent(\pi) = -\sum_a \pi(a)\ln \pi(a) ,
\]
which is $0$ when one move takes everything and $\ln K$ when $K$ moves share equally. For the
drawn ending, $\pi = (0.216, 0.214, 0.194, 0.190, 0.185)$:
\[
\Ent = -[\,0.216\ln 0.216 + 0.214\ln 0.214 + 0.194\ln 0.194 + 0.190\ln 0.190 + 0.185\ln 0.185\,]
= 1.607,
\]
against a maximum of $\ln 5 = 1.609$. The search spread its budget as evenly as it is
possible to spread it. Here is the same measure across all five positions in this book's
data set:
\begin{center}
\begin{tabular}{@{}lrrrl@{}}
\toprule
Position & moves visited & $\Ent(\pi)$ & $\ln(\text{legal})$ & shape \\
\midrule
Drawn K+P ending      & 5 of 5   & $1.607$ & $1.609$ & flat (nothing at stake) \\
Initial position      & 20 of 20 & $1.701$ & $2.996$ & split (several playable) \\
Won K+P ending        & 4 of 4   & $0.690$ & $1.386$ & concentrated on two \\
Morphy, before O-O-O  & 9 of 49  & $0.490$ & $3.892$ & one move \\
Morphy, before Qb8+   & 8 of 46  & $0.427$ & $3.829$ & one move \\
\bottomrule
\end{tabular}
\end{center}
Note the fourth column: in the Morphy positions the search visited 8 or 9 of the $\approx 47$
legal moves \emph{at all}. Roughly four out of five legal moves received zero attention. That
is not a defect --- it is requirement R1 of Chapter~\ref{ch:puct}, working.

Entropy is the same tool that will define "sharpness" for an \emph{opponent's} options in
Chapter~\ref{ch:style}; it is built from scratch there. Here it is just a shape summary.
\end{worked}

\section{How much did the search add?}

The network's static opinion of a move is $V$ --- one forward pass, zero search. The
post-search opinion is $Q$. The difference
\begin{equation}
  \label{eq:searchgain}
  \Delta_{\text{search}}(a) \;=\; Q(s,a) - V(s\!\cdot\!a)
\end{equation}
is a directly useful quantity: \emph{how much did calculation change the assessment of this
move?} And unlike most things in this book, it is printed for you.

\begin{realdata}[title={Real engine output: what search added, Morphy position}]
\begin{center}
\begin{tabular}{@{}lrrrr@{}}
\toprule
Move & $n$ & $V$ (no search) & $Q$ (after search) & $\Delta_{\text{search}}$ \\
\midrule
Qb7   & 172 & $+0.3045$ & $+0.6587$ & $\mathbf{+0.354}$ \\
Rxd7  &   4 & $+0.1415$ & $+0.3155$ & $+0.174$ \\
Qb5   &   2 & $-0.1652$ & $-0.1019$ & $+0.063$ \\
Qxe6+ &   1 & $-0.3595$ & $-0.3570$ & $+0.003$ \\
Qb8+  &   3 & $+1.0000$ & $+1.0000$ & $\ \ 0.000$ (proven) \\
\bottomrule
\end{tabular}
\end{center}
The network looked at the position after Qb7 and said "$+0.30$: better". After 172 visits the
search said "$+0.66$: much better". \textbf{Calculation was worth $0.35$ of evaluation here}
--- and note it is worth more the more it was applied: the four-visit move gained $0.17$, the
one-visit move gained nothing, because a single visit is just the static evaluation again.
\end{realdata}

\begin{keyidea}
$\Delta_{\text{search}}$ is a measure of \textbf{how much a position rewards calculation}. It
is large where the static picture is misleading and small where the position is what it looks
like. For a training tool this is directly actionable: positions with large
$\Delta_{\text{search}}$ across many moves are positions where \emph{your} intuition is also
likely to be insufficient --- the ones worth putting on a clock and working out.
\end{keyidea}

\section{Policy versus visits: the divergence index}

Now the comparison your project is built on. Define the visit share
$\pi(a) = N(s,a)/N(s)$ and compare it with the prior:
\begin{equation}
  \label{eq:divergence}
  \Delta_{\text{P,MCTS}}(a) \;=\; \pi(a) \;-\; P(a\given s) .
\end{equation}
By Chapter~\ref{ch:puct}'s equilibrium theorem, this is exactly \textbf{the information the
search added to the network's instinct}. Positive: the move earned more attention than its
instinct deserved. Negative: it was demoted.

\begin{center}
\begin{tabular}{@{}llll@{}}
\toprule
 & $P$ high & $P$ low \\
\midrule
$\pi$ high & \textbf{agreement} --- natural and sound & \textbf{hidden gem} --- counter-intuitive and strong \\
$\pi$ low  & \textbf{optical trap} --- natural and refuted & (ignored, correctly) \\
\bottomrule
\end{tabular}
\end{center}

Two of these categories are worth a great deal to a training tool, for opposite reasons. An
\emph{optical trap} is a move your own instinct is likely to share --- the network's policy
is a model of what a strong player plays, so a move it likes is a move that looks right. A
\emph{hidden gem} is the opposite: a move that would not have occurred to a strong player,
and did not occur to the network, but survives calculation.

Here are all three, in real data.

\begin{realdata}[title={Real engine output: three positions, three verdicts}]
\textbf{Agreement} --- Morphy position, four moves earlier (before 12.O-O-O), 800 nodes:
\begin{center}
\begin{tabular}{@{}lrrr@{}}
\toprule
Move & $P$ & $\pi$ & $\Delta_{\text{P,MCTS}}$ \\
\midrule
O-O-O & 35.6\% & 87.3\% & $+51.7$ pts \\
Bxf6  & 21.2\% & 10.1\% & $-11.1$ pts \\
Rd1   &  7.1\% &  0.8\% & $-6.3$ pts \\
\bottomrule
\end{tabular}
\end{center}
The network's top choice was also the search's, reinforced: instinct and calculation agree.
(Stockfish at depth 30 concurs: O-O-O is best at $+5.19$, more than double the next move.)

\medskip
\textbf{Optical trap (small)} --- the K+P endgame of Chapter~\ref{ch:byhand}, 800 nodes:
\begin{center}
\begin{tabular}{@{}lrrrl@{}}
\toprule
Move & $P$ & $\pi$ & $\Delta_{\text{P,MCTS}}$ & \\
\midrule
Kd6 & 45.1\% & 60.9\% & $+15.8$ pts & wins (mate in 12) \\
Kf6 & 44.2\% & 38.8\% & $-5.4$ pts & wins (mate in 18) \\
Kf5 &  5.4\% &  0.2\% & $-5.2$ pts & \textbf{draws} \\
Kd5 &  5.3\% &  0.2\% & $-5.1$ pts & \textbf{draws} \\
\bottomrule
\end{tabular}
\end{center}
The instinct assigned $10.6\%$ to two moves that throw away a won game; search demoted them
to $0.4\%$ between them. Small in absolute terms, but it is the same mechanism as a
game-losing trap, and the ground truth (Stockfish, depth 30) confirms the demotion was right.

\medskip
\textbf{Hidden gem} --- Morphy position before 16.Qb8+:
\begin{center}
\begin{tabular}{@{}lrrrl@{}}
\toprule
Move & $P$ & $\pi$ & $\Delta_{\text{P,MCTS}}$ & \\
\midrule
Qb7  & 31.9\% & 92.0\% & $+60.1$ pts & good; not best \\
\textbf{Qb8+} & \textbf{1.6\%} & 1.6\% & $\approx 0$ & \textbf{mate in 2 --- and the move played} \\
\bottomrule
\end{tabular}
\end{center}
\end{realdata}

\begin{pitfall}
\textbf{Look carefully at that last table, because it breaks the neat 2$\times$2.} The gem was
found --- the engine plays Qb8+ --- but its \emph{visit share is 1.6\%}, identical to its
prior. The divergence index is about zero. If your tool classified positions by
Equation~\eqref{eq:divergence} alone, it would score the most famous queen sacrifice in chess
literature as "nothing happened here".

Why? Because the move was proven, not sampled (Chapter~\ref{ch:engineering}). Proof stops the
search instead of accumulating visits. \textbf{The divergence index measures attention, and
proof is cheap in attention.} A usable classifier must therefore look at three things
together: divergence, $Q$, and whether $Q$ is certain. Something like:
\[
\text{gem}(a) \;=\; \bigl[\,P(a\given s) \text{ low}\,\bigr] \;\wedge\;
\bigl[\,Q(s,a) \text{ best at the root}\,\bigr] \;\wedge\;
\bigl[\,\pi(a) \text{ high \emph{or} } Q(s,a) \text{ proven}\,\bigr].
\]
This is a real correction to the design as it stands in your notes, and it came from looking
at one position's actual numbers rather than from theory.
\end{pitfall}

\section{How many decimals may you quote?}

From Chapter~\ref{ch:onemachine}: the standard error of a mean over $n$ samples is
$\sigma/\sqrt{n}$. Taking $\sigma \approx 0.4$ for values on $[-1,1]$, here is the table to
keep beside your display code:

\begin{center}
\begin{tabular}{@{}rll@{}}
\toprule
$n$ & typical error & honest presentation \\
\midrule
1--3     & $\pm 0.2$--$0.4$ & "the network's guess", no number \\
10       & $\pm 0.13$ & one decimal at most, hedged \\
100      & $\pm 0.04$ & one decimal \\
1{,}000  & $\pm 0.013$ & two decimals \\
10{,}000 & $\pm 0.004$ & two decimals \\
\bottomrule
\end{tabular}
\end{center}

\begin{keyidea}
A move with $n = 3$ has no meaningful evaluation, even though the engine prints five decimal
places for it. The commonest quiet dishonesty in chess software is displaying a
well-formatted number that the search did not earn. \textbf{Attach $n$ to every evaluation
your tool shows, or suppress the evaluation.} The exception, again, is a proven value.
\end{keyidea}

\section{A reading checklist}

Put this next to the screen.

\begin{enumerate}
\item \textbf{Root visits?} Below a few hundred, treat everything as provisional.
\item \textbf{Draw probability?} High $d$ means the rest of the analysis is about very little.
\item \textbf{Visit shape?} Concentrated, split, or flat --- and if flat, is that the position
  or the budget?
\item \textbf{Top move's $n$ and $Q$.} Quote decimals per the table above.
\item \textbf{$Q$ versus $V$.} Large gap = the position rewards calculation.
\item \textbf{$\pi$ versus $P$.} Large positive on a low-prior move = look here. Large
  negative on a high-prior move = a trap your instinct probably shares.
\item \textbf{Any $Q = \pm1.000$ at small $n$?} That is a proof; treat it as decisive and
  ignore its visit count.
\item \textbf{PV last}, and only as far as the visit counts along it support.
\end{enumerate}

\begin{repobox}[title={in this repository}]
Items 5, 6 and 7 are the raw material of the explanation layer: $\Delta_{\text{search}}$ says
\emph{whether} there is something to explain, $\Delta_{\text{P,MCTS}}$ says \emph{which move}
the explanation is about, and the certainty flag says \emph{how strongly} to phrase it.
Chapter~\ref{ch:speak} turns these into sentences; Chapter~\ref{ch:capstone} runs one
position all the way through.
\end{repobox}

\begin{exercises}

\exhead[reading order]{ch9:order} You are handed a dump with root $N = 47$ and a beautiful
12-ply PV. Which items of §9.1 can you use, and what should you refuse to say?

\exhead[entropy]{ch9:entropy} Compute the entropy of the visit distributions
$(0.90,0.04,0.03,0.02,0.01)$ and $(0.30,0.28,0.22,0.12,0.08)$. Which is which shape from
§9.2, and what is the maximum possible entropy for five moves?

\exhead[search gain]{ch9:gain} From the real table in §9.3, why does Qxe6+ have
$\Delta_{\text{search}} \approx 0$? Is that evidence that the move is well understood?
Explain carefully --- this is a trick question about $n$.

\exhead[divergence arithmetic]{ch9:divarith} A move has $P = 4\%$ and $N(s,a) = 300$ out of
$N(s) = 1000$. Compute $\pi$ and $\Delta_{\text{P,MCTS}}$. Classify the move. What single
extra number would you want before showing this to a user as a "hidden gem"?

\exhead[the trap]{ch9:trap} A move has $P = 62\%$ and $\pi = 3\%$. Write the sentence your
tool should show, using no chess content beyond the move's name --- and say what makes this
category especially valuable to a player whose diagnosed weakness is tactics.

\exhead[the gem that scored zero]{ch9:gemzero} Restate, in your own words, why the Qb8+
divergence was $\approx 0$ despite the move being a forced mate. Then propose a second
detector that would have caught it, and give the position from this book on which you would
test it.

\exhead[decimals]{ch9:decimals} Your tool shows "$-0.13$" for a move with $n = 6$. Using the
error table, what is the honest statement? Rewrite the display text.

\exhead[a fake gem]{ch9:fakegem} Construct a plausible set of numbers ($P$, $n$, $Q$, $V$) for
a move that would pass a naive "low policy, high $Q$" gem detector and yet be nothing of the
sort. What check kills it?

\exhead[flat is ambiguous]{ch9:flat} Two positions both show a flat visit distribution. One
has root $N = 60$, the other $N = 4000$ and $d = 0.99$. Write the different sentence your
tool should produce for each.

\exhead[the checklist in practice]{ch9:practice} Take any position from your own games,
analyse it, and run the eight-item checklist end to end in writing. Then write the two
sentences you would want the tool to have said. Keep them; Chapter~\ref{ch:speak} asks you to
compare.

\end{exercises}
